\section{Conclusão}
\label{sec:conclusao}

Nesta primeira parte, o problema foi formulado, os dados foram limpos e explorados e foi definido o plano de modelação, sempre com o objetivo de evitar fugas de informação. As decisões estruturais foram tomadas antes da análise exploratória: as tarefas e os seus momentos de previsão (antes do \textit{pricing} na regressão, na concessão na classificação), a elegibilidade de cada uma das 151 colunas, as coortes e uma divisão temporal com teste em 2015. A exigência de que os rótulos de treino fossem observáveis na data simulada tornou a classificação, na prática, binária: a classe \textit{At risk} não existe no treino e, no teste, é contada como \textit{Default}.

A análise exploratória, feita apenas com dados de desenvolvimento, mostrou que a taxa de juro depende sobretudo da qualidade do crédito e do prazo, mas que nenhuma variável isolada a explica; que o grau atribuído pela plataforma separa bem o risco de \textit{default} (de 5,9\% a 35,9\%), enquanto cada variável numérica, sozinha, separa pouco as classes; que várias relações com o \textit{default} não são lineares; que há variáveis quase redundantes e caudas longas, que justificam o logaritmo e a limitação dos valores extremos; e que a população mudou entre 2007--2011 e 2013--2014. Os testes de hipóteses, feitos em R com a correção de Holm, rejeitaram a hipótese nula nos quatro casos, mas com efeitos muito diferentes: desprezável no DTI ($r_{rb} = 0{,}065$), pequeno na finalidade ($V = 0{,}104$) e no montante pedido ($r_{rb} = 0{,}197$) e grande apenas no \textit{risk score}, mais baixo nos pedidos rejeitados ($r_{rb} = -0{,}614$).

Com base nestas conclusões, foram criadas dez \textit{features}, cada uma com a elegibilidade herdada das colunas de origem, e foi definido um plano com kNN implementado de raiz, modelos lineares regularizados, árvores de decisão, \textit{Random Forest}, LightGBM e uma rede neuronal densa, validados com uma janela temporal expansível e comparados com \textit{baselines} simples.

As principais limitações são três. Primeiro, só se observam empréstimos aceites, pelo que os modelos não devem ser aplicados a pedidos fora das regras de concessão da plataforma. Segundo, o classificador é treinado com empréstimos de 2007--2011 e avaliado em 2015, numa população com pior perfil de crédito, pelo que o desempenho no teste pode ser inferior ao da validação. Terceiro, não é possível identificar mutuários repetidos e assume-se que os campos do \textit{bureau} correspondem ao momento do pedido.

Na Parte~2, os modelos serão implementados e afinados com os mesmos \textit{folds} e avaliados uma única vez no teste de 2015, com as métricas e \textit{baselines} aqui definidas; serão construídos os agrupamentos de mutuários e de estados, avaliando a sua estabilidade; será analisada a utilização das variáveis geográficas do ponto de vista da equidade (\textit{fairness}), que por agora ficaram de fora; e será descrito o plano de implementação do modelo final.